% GENERATED FILE. Edit canonical lessons, then run npm run generate:books.
% canonical-source-hash: fnv1a64:f2df8ed38e9957d2
% canonical-lessons: LA-C305-scelus, LA-C305-facinus, LA-C305-sceleratus, LA-C305-fraus, LA-C305-dolus

\chapter{Scelus et fraus}
\label{ch:la-scelus-et-fraus}
\hlchaptermodality{305}

\begin{chapteropening}
\textbf{By the end of this chapter:} \emph{I can name wrongdoing in a story: a crime, a deed, wicked, fraud and trickery.}

Complete the last lesson of chapter 305: I can name wrongdoing in a story: a crime, a deed, wicked, fraud and trickery.
\end{chapteropening}

\section[scelus, sceleris]{scelus, sceleris — a crime, a wicked deed}

\label{lesson:LA-C305-scelus}



\begin{quote}
Before the new one: say the Latin for to catch, then the Latin for to seize.
\end{quote}

\subsection*{You'll want to know: scelus, sceleris}
\textbf{scelus, sceleris} — \textquotedblleft{}a crime, a wicked deed\textquotedblright{}.

\begin{quote}
Omnēs cīvēs dē hōc scelere loquēbantur.
\end{quote}

\begin{quote}
\emph{All the citizens were talking about this crime.}
\end{quote}

Neuter, 3rd declension. The strongest word for a wrong: a crime that also offends the gods.

\begin{cousinweb}
\textbf{scelus} is an inherited word whose origin is not settled; Lewis \& Short compare a Sanskrit root meaning \textquotedblleft{}to stumble\textquotedblright{}.
\end{cousinweb}

\begin{grammarlens}[title={What you've built}]
One more word for telling, explaining and describing at B1.
\end{grammarlens}

\subsection*{Your turn}
\begin{itemize}
\raggedright
  \item \emph{Say it:} \emph{scelus, sceleris}
  \item \emph{Say it:} \emph{scelus, sceleris}, once more
  \item \emph{Say it:} a sentence of your own with \emph{scelus, sceleris}
  \item \emph{Recall:} say the Latin for to catch, then the Latin for to seize, then say \emph{scelus, sceleris} again
\end{itemize}

\begin{tcolorbox}[breakable,colback=teal!4,colframe=teal!35!black,title={Before you move on}]
What does scelus, sceleris mean? (\textquotedblleft{}A crime, a wicked deed\textquotedblright{}.) Say it once more.
\end{tcolorbox}
\section[facinus, facinoris]{facinus, facinoris — a deed; a bad deed}

\label{lesson:LA-C305-facinus}



\begin{quote}
Before the new one: say the Latin for to seize, then the Latin for a crime.
\end{quote}

\subsection*{You'll want to know: facinus, facinoris}
\textbf{facinus, facinoris} — \textquotedblleft{}a deed; a bad deed\textquotedblright{}.

\begin{quote}
Facinus audāx erat, sed nōn turpe.
\end{quote}

\begin{quote}
\emph{It was a daring deed, but not a shameful one.}
\end{quote}

Neuter, 3rd declension, from \emph{faciō}. Alone it often means a crime; praised, a fine deed.

\begin{cousinweb}
\textbf{facinus} is built on \emph{faciō}, \textquotedblleft{}to do, to make\textquotedblright{} (EtymDB): a thing done, and very often a wicked thing done.
\end{cousinweb}

\begin{grammarlens}[title={What you've built}]
One more word for telling, explaining and describing at B1.
\end{grammarlens}

\subsection*{Your turn}
\begin{itemize}
\raggedright
  \item \emph{Say it:} \emph{facinus, facinoris}
  \item \emph{Say it:} \emph{facinus, facinoris}, once more
  \item \emph{Say it:} a sentence of your own with \emph{facinus, facinoris}
  \item \emph{Recall:} say the Latin for to seize, then the Latin for a crime, then say \emph{facinus, facinoris} again
\end{itemize}

\begin{tcolorbox}[breakable,colback=teal!4,colframe=teal!35!black,title={Before you move on}]
What does facinus, facinoris mean? (\textquotedblleft{}A deed; a bad deed\textquotedblright{}.) Say it once more.
\end{tcolorbox}
\section[scelerātus, scelerāta, scelerātum]{scelerātus, scelerāta, scelerātum — wicked, criminal}

\label{lesson:LA-C305-sceleratus}



\begin{quote}
Before the new one: say the Latin for a crime, then the Latin for a deed.
\end{quote}

\subsection*{You'll want to know: scelerātus, scelerāta, scelerātum}
\textbf{scelerātus, scelerāta, scelerātum} — \textquotedblleft{}wicked, criminal\textquotedblright{}.

\begin{quote}
Ille vir scelerātus patrem suum fefellit.
\end{quote}

\begin{quote}
\emph{That wicked man deceived his own father.}
\end{quote}

The past participle of \emph{scelerō}, 'to stain with crime', used as an adjective.

\begin{cousinweb}
\textbf{scelerātus} is the participle of \emph{scelerō}, \textquotedblleft{}to pollute with crime\textquotedblright{}, built on \emph{scelus}, \textquotedblleft{}crime\textquotedblright{} (Lewis \& Short).
\end{cousinweb}

\begin{grammarlens}[title={What you've built}]
One more word for telling, explaining and describing at B1.
\end{grammarlens}

\subsection*{Your turn}
\begin{itemize}
\raggedright
  \item \emph{Say it:} \emph{scelerātus, scelerāta, scelerātum}
  \item \emph{Say it:} \emph{scelerātus, scelerāta, scelerātum}, once more
  \item \emph{Say it:} a sentence of your own with \emph{scelerātus, scelerāta, scelerātum}
  \item \emph{Recall:} say the Latin for a crime, then the Latin for a deed, then say \emph{scelerātus, scelerāta, scelerātum} again
\end{itemize}

\begin{tcolorbox}[breakable,colback=teal!4,colframe=teal!35!black,title={Before you move on}]
What does scelerātus, scelerāta, scelerātum mean? (\textquotedblleft{}Wicked, criminal\textquotedblright{}.) Say it once more.
\end{tcolorbox}
\section[fraus, fraudis]{fraus, fraudis — fraud, deceit; harm}

\label{lesson:LA-C305-fraus}



\begin{quote}
Before the new one: say the Latin for a deed, then the Latin for wicked.
\end{quote}

\subsection*{You'll want to know: fraus, fraudis}
\textbf{fraus, fraudis} — \textquotedblleft{}fraud, deceit; harm\textquotedblright{}.

\begin{quote}
Mercātor fraude multam pecūniam accēpit.
\end{quote}

\begin{quote}
\emph{The merchant got a lot of money by fraud.}
\end{quote}

Feminine, 3rd declension. In old law also 'harm': \emph{sine fraude}, without loss. English \emph{fraud}.

\begin{cousinweb}
\textbf{fraus} is an inherited word whose earlier origin is not settled (Lewis \& Short, EtymDB). English \emph{fraud} comes from it through French.
\end{cousinweb}

\begin{grammarlens}[title={What you've built}]
One more word for telling, explaining and describing at B1.
\end{grammarlens}

\subsection*{Your turn}
\begin{itemize}
\raggedright
  \item \emph{Say it:} \emph{fraus, fraudis}
  \item \emph{Say it:} \emph{fraus, fraudis}, once more
  \item \emph{Say it:} a sentence of your own with \emph{fraus, fraudis}
  \item \emph{Recall:} say the Latin for a deed, then the Latin for wicked, then say \emph{fraus, fraudis} again
\end{itemize}

\begin{tcolorbox}[breakable,colback=teal!4,colframe=teal!35!black,title={Before you move on}]
What does fraus, fraudis mean? (\textquotedblleft{}Fraud, deceit; harm\textquotedblright{}.) Say it once more.
\end{tcolorbox}
\section[dolus, dolī]{dolus, dolī — a trick, cunning}

\label{lesson:LA-C305-dolus}



\begin{quote}
Before the new one: say the Latin for wicked, then the Latin for fraud.
\end{quote}

\subsection*{You'll want to know: dolus, dolī}
\textbf{dolus, dolī} — \textquotedblleft{}a trick, cunning\textquotedblright{}.

\begin{quote}
Hostēs dolum parāvērunt, sed dux eum animadvertit.
\end{quote}

\begin{quote}
\emph{The enemy prepared a trick, but the leader noticed it.}
\end{quote}

Masculine, 2nd declension. The wooden horse was the most famous \emph{dolus} of the old stories.

\begin{cousinweb}
\textbf{dolus} matches Greek \emph{dolos}, \textquotedblleft{}bait, trick\textquotedblright{}; EtymDB lists it as taken from Greek, while Lewis \& Short treat the two as cousins.
\end{cousinweb}

\begin{grammarlens}[title={What you've built}]
One more word for telling, explaining and describing at B1.
\end{grammarlens}

\subsection*{Your turn}
\begin{itemize}
\raggedright
  \item \emph{Say it:} \emph{dolus, dolī}
  \item \emph{Say it:} \emph{dolus, dolī}, once more
  \item \emph{Say it:} a sentence of your own with \emph{dolus, dolī}
  \item \emph{Recall:} say the Latin for wicked, then the Latin for fraud, then say \emph{dolus, dolī} again
\end{itemize}

\begin{tcolorbox}[breakable,colback=teal!4,colframe=teal!35!black,title={Before you move on}]
What does dolus, dolī mean? (\textquotedblleft{}A trick, cunning\textquotedblright{}.) Say it once more.
\end{tcolorbox}
